\begin{remark*}[Exposition]
Tao introduces finite series by first fixing a finite real sequence indexed
by consecutive integers and then defining its sum recursively
\cite{TaoAnalysis1}. The source bundles these conventions into one
definition; here they are split into atomic definitions.
\end{remark*}

\begin{definitionbox}{Definition (Finite Real Sequence on an Integer Interval)}
% ---------------------------------------------------------
% Definition: Finite Real Sequence on an Integer Interval
% Label: def:finite-real-sequence-integer-interval
% Source: Tao, Analysis I, Definition 7.1.1
% ---------------------------------------------------------
\begin{definition}[Finite Real Sequence on an Integer Interval]
\label{def:finite-real-sequence-integer-interval}

Let \(m,n \in \mathbb{Z}\) with \(m \le n\). A finite real sequence indexed
from \(m\) to \(n\) is an assignment of a real number \(a_i\) to each
integer \(i\) satisfying \(m \le i \le n\). We denote such a sequence by
\((a_i)_{i=m}^{n}\).
\end{definition}
\end{definitionbox}

\begin{remark*}[Standard quantified statement]
\[
  m,n\in\mathbb Z\wedge m\le n
  \Longrightarrow
  (a_i)_{i=m}^{n}\text{ assigns a real number }a_i
  \text{ to each }i\in\mathbb Z\text{ with }m\le i\le n.
\]
\end{remark*}

\begin{remark*}[Interpretation]
The indices form the integer interval \(\{m,m+1,\ldots,n\}\). The notation
\((a_i)_{i=m}^{n}\) records both the terms and the finite range of indices
on which the terms are defined.
\end{remark*}

\begin{dependencies}
\begin{itemize}
  \item \hyperref[def:real-sequence]{Real Sequence}
\end{itemize}
\end{dependencies}

\begin{definitionbox}{Definition (Empty Finite Sum)}
% ---------------------------------------------------------
% Definition: Empty Finite Sum
% Label: def:empty-finite-sum
% Source: Tao, Analysis I, Definition 7.1.1
% ---------------------------------------------------------
\begin{definition}[Empty Finite Sum]
\label{def:empty-finite-sum}

Let \(m,n \in \mathbb{Z}\). If \(n < m\), the finite sum from \(m\) to
\(n\) is defined by
\[
  \sum_{i=m}^{n} a_i := 0.
\]
\end{definition}
\end{definitionbox}

\begin{remark*}[Standard quantified statement]
\[
  n<m\Longrightarrow \sum_{i=m}^{n}a_i=0.
\]
\end{remark*}

\begin{remark*}[Interpretation]
When the upper index lies below the lower index, there are no terms to
add. The sum is therefore assigned the additive identity \(0\).
\end{remark*}

\begin{dependencies}
\begin{itemize}
  \item \hyperref[def:finite-real-sequence-integer-interval]{Finite Real Sequence on an Integer Interval}
\end{itemize}
\end{dependencies}

\begin{definitionbox}{Definition (Finite Sum)}
% ---------------------------------------------------------
% Definition: Finite Sum
% Label: def:finite-sum
% Source: Tao, Analysis I, Definition 7.1.1
% ---------------------------------------------------------
\begin{definition}[Finite Sum]
\label{def:finite-sum}

Let \(m,n \in \mathbb{Z}\), and let real numbers \(a_i\) be assigned
for the relevant integer indices. The notation
\[
  \sum_{i=m}^{n} a_i
\]
denotes the finite sum, or finite series, of the terms indexed from
\(m\) to \(n\).
\end{definition}
\end{definitionbox}

\begin{remark*}[Standard quantified statement]
\[
  \sum_{i=m}^{n} a_i
  \quad\text{denotes the finite sum of the terms } a_i
  \text{ indexed by } m \le i \le n.
\]
\end{remark*}

\begin{remark*}[Interpretation]
This definition fixes the notation and name of the object being defined.
The value of the object is then determined by the empty-sum convention
and the recursive rule below.
\end{remark*}

\begin{dependencies}
\begin{itemize}
  \item \hyperref[def:finite-real-sequence-integer-interval]{Finite Real Sequence on an Integer Interval}
\end{itemize}
\end{dependencies}

\begin{definitionbox}{Definition (Recursive Finite Sum)}
% ---------------------------------------------------------
% Definition: Recursive Finite Sum
% Label: def:recursive-finite-sum
% Source: Tao, Analysis I, Definition 7.1.1
% ---------------------------------------------------------
\begin{definition}[Recursive Finite Sum]
\label{def:recursive-finite-sum}

Let \(m \in \mathbb{Z}\), and let real numbers \(a_i\) be assigned for
the relevant indices. The finite sum, also called a finite series, is
defined recursively from the empty finite sum by the rule
\[
  \sum_{i=m}^{n+1} a_i
  :=
  \left(\sum_{i=m}^{n} a_i\right) + a_{n+1}
  \qquad \text{whenever } n \ge m-1.
\]
\end{definition}
\end{definitionbox}

\begin{remark*}[Standard quantified statement]
\[
  \forall m,n \in \mathbb{Z},\quad
  n \ge m-1
  \Longrightarrow
  \sum_{i=m}^{n+1} a_i
  :=
  \left(\sum_{i=m}^{n} a_i\right)+a_{n+1}.
\]
\end{remark*}

\begin{remark*}[Predicate reading]
\[
  n\ge m-1
  \Longrightarrow
  \sum_{i=m}^{n+1}a_i
  =
  \left(\sum_{i=m}^{n}a_i\right)+a_{n+1}.
\]
\end{remark*}

\begin{remark*}[Interpretation]
The recursive rule says that extending the upper endpoint by one index
extends the sum by adding exactly the new last term. Together with the
empty-sum convention, this determines every finite sum.
\end{remark*}

\begin{remark*}[Interpretation]
For a sequence indexed from \(m\), the first few cases are
\[
  \sum_{i=m}^{m-2} a_i = 0,\qquad
  \sum_{i=m}^{m-1} a_i = 0,\qquad
  \sum_{i=m}^{m} a_i = a_m,
\]
and
\[
  \sum_{i=m}^{m+1} a_i = a_m+a_{m+1},\qquad
  \sum_{i=m}^{m+2} a_i = a_m+a_{m+1}+a_{m+2}.
\]
\end{remark*}

\begin{dependencies}
\begin{itemize}
  \item \hyperref[def:empty-finite-sum]{Empty Finite Sum}
  \item \hyperref[def:finite-sum]{Finite Sum}
  \item \hyperref[def:finite-real-sequence-integer-interval]{Finite Real Sequence on an Integer Interval}
\end{itemize}
\end{dependencies}

\begin{lemmabox}{Lemma (Splitting a Finite Sum)}
% ---------------------------------------------------------
% Lemma: Splitting a Finite Sum
% Label: lem:finite-sum-splitting
% Source: Tao, Analysis I, Lemma 7.1.4(a)
% ---------------------------------------------------------
\begin{lemma}[Splitting a Finite Sum]
\label{lem:finite-sum-splitting}

Let \(m \le n < p\) be integers, and let \(a_i\) be a real number
assigned to each integer \(i\) satisfying \(m \le i \le p\). Then
\[
  \sum_{i=m}^{n} a_i
  +
  \sum_{i=n+1}^{p} a_i
  =
  \sum_{i=m}^{p} a_i.
\]

\smallskip
\noindent
\hyperref[prf:finite-sum-splitting]{\textit{Go to proof.}}
\end{lemma}
\end{lemmabox}

\begin{remark*}[Standard quantified statement]
\[
  m\le n<p
  \Longrightarrow
  \sum_{i=m}^{n}a_i+\sum_{i=n+1}^{p}a_i
  =
  \sum_{i=m}^{p}a_i.
\]
\end{remark*}

\begin{remark*}[Interpretation]
A finite sum over one integer interval may be split into two adjacent
finite sums without changing its value.
\end{remark*}

\begin{dependencies}
\begin{itemize}
  \item \hyperref[def:recursive-finite-sum]{Recursive Finite Sum}
\end{itemize}
\end{dependencies}

\begin{definitionbox}{Definition (Summation over a Finite Set)}
% ---------------------------------------------------------
% Definition: Summation over a Finite Set
% Label: def:summation-over-finite-set
% Source: Tao, Analysis I, Definition 7.1.6
% ---------------------------------------------------------
\begin{definition}[Summation over a Finite Set]
\label{def:summation-over-finite-set}

Let \(X\) be a finite set with \(n\) elements, where \(n \in \mathbb{N}\),
and let \(f:X\to\mathbb{R}\). Choose a bijection
\[
  g:\{i\in\mathbb{N}:1\le i\le n\}\to X.
\]
The finite sum of \(f\) over \(X\) is defined by
\[
  \sum_{x\in X} f(x)
  :=
  \sum_{i=1}^{n} f(g(i)).
\]
\end{definition}
\end{definitionbox}

\begin{remark*}[Standard quantified statement]
\[
  X=\{g(1),\ldots,g(n)\}
  \quad\Longrightarrow\quad
  \sum_{x\in X} f(x)
  :=
  \sum_{i=1}^{n} f(g(i)).
\]
\end{remark*}

\begin{remark*}[Interpretation]
The definition computes a sum over an unordered finite set by first listing
its elements through a bijection from \(\{1,\ldots,n\}\). The notation
\(\sum_{x\in X}f(x)\) records the set being summed over; the chosen listing
is auxiliary.
\end{remark*}

\begin{dependencies}
\begin{itemize}
  \item \hyperref[def:finite-sum]{Finite Sum}
\end{itemize}
\end{dependencies}

\begin{propositionbox}{Proposition (Finite Summations Are Well-Defined)}
% ---------------------------------------------------------
% Proposition: Finite Summations Are Well-Defined
% Label: prop:finite-summations-well-defined
% Source: Tao, Analysis I, Proposition 7.1.8
% ---------------------------------------------------------
\begin{proposition}[Finite Summations Are Well-Defined]
\label{prop:finite-summations-well-defined}

Let \(X\) be a finite set with \(n\) elements, where \(n \in \mathbb{N}\),
let \(f:X\to\mathbb{R}\), and let
\[
  g:\{i\in\mathbb{N}:1\le i\le n\}\to X
  \qquad\text{and}\qquad
  h:\{i\in\mathbb{N}:1\le i\le n\}\to X
\]
be bijections. Then
\[
  \sum_{i=1}^{n} f(g(i))
  =
  \sum_{i=1}^{n} f(h(i)).
\]

\smallskip
\noindent
\hyperref[prf:finite-summations-well-defined]{\textit{Go to proof.}}
\end{proposition}
\end{propositionbox}

\begin{remark*}[Standard quantified statement]
\[
  \sum_{i=1}^{n}f(g(i))
  =
  \sum_{i=1}^{n}f(h(i))
\]
whenever \(g,h:\{i\in\mathbb N:1\le i\le n\}\to X\) are bijections.
\end{remark*}

\begin{remark*}[Interpretation]
The value of a finite summation over a finite set does not depend on which
bijection is used to list the elements of that set. Tao notes that the
corresponding issue for summations over infinite sets is more complicated.
\end{remark*}

\begin{dependencies}
\begin{itemize}
  \item \hyperref[def:summation-over-finite-set]{Summation over a Finite Set}
\end{itemize}
\end{dependencies}

\begin{propositionbox}{Proposition (Empty Finite Set Summation)}
% ---------------------------------------------------------
% Proposition: Empty Finite Set Summation
% Label: prop:empty-finite-set-summation
% Source: Tao, Analysis I, Proposition 7.1.11(a)
% ---------------------------------------------------------
\begin{proposition}[Empty Finite Set Summation]
\label{prop:empty-finite-set-summation}

If \(X\) is empty, and \(f:X\to\mathbb{R}\) is a function, then
\[
  \sum_{x\in X} f(x)=0.
\]

\smallskip
\noindent
\hyperref[prf:empty-finite-set-summation]{\textit{Go to proof.}}
\end{proposition}
\end{propositionbox}

\begin{remark*}[Standard quantified statement]
\[
  X=\emptyset\Longrightarrow \sum_{x\in X}f(x)=0.
\]
\end{remark*}

\begin{remark*}[Interpretation]
The finite summation over the empty set has value \(0\).
\end{remark*}

\begin{dependencies}
\begin{itemize}
  \item \hyperref[def:summation-over-finite-set]{Summation over a Finite Set}
\end{itemize}
\end{dependencies}

\begin{propositionbox}{Proposition (Singleton Finite Set Summation)}
% ---------------------------------------------------------
% Proposition: Singleton Finite Set Summation
% Label: prop:singleton-finite-set-summation
% Source: Tao, Analysis I, Proposition 7.1.11(b)
% ---------------------------------------------------------
\begin{proposition}[Singleton Finite Set Summation]
\label{prop:singleton-finite-set-summation}

If \(X\) consists of a single element, \(X=\{x_0\}\), and
\(f:X\to\mathbb{R}\) is a function, then
\[
  \sum_{x\in X} f(x)=f(x_0).
\]

\smallskip
\noindent
\hyperref[prf:singleton-finite-set-summation]{\textit{Go to proof.}}
\end{proposition}
\end{propositionbox}

\begin{remark*}[Standard quantified statement]
\[
  X=\{x_0\}\Longrightarrow \sum_{x\in X}f(x)=f(x_0).
\]
\end{remark*}

\begin{remark*}[Interpretation]
Summing a function over a one-element set gives the value of the function
at that element.
\end{remark*}

\begin{dependencies}
\begin{itemize}
  \item \hyperref[def:summation-over-finite-set]{Summation over a Finite Set}
\end{itemize}
\end{dependencies}

\begin{propositionbox}{Proposition (Substitution for Finite Set Summation)}
% ---------------------------------------------------------
% Proposition: Substitution for Finite Set Summation
% Label: prop:substitution-finite-set-summation
% Source: Tao, Analysis I, Proposition 7.1.11(c)
% ---------------------------------------------------------
\begin{proposition}[Substitution for Finite Set Summation]
\label{prop:substitution-finite-set-summation}

If \(X\) is a finite set, \(f:X\to\mathbb{R}\) is a function, and
\(g:Y\to X\) is a bijection, then
\[
  \sum_{x\in X} f(x)
  =
  \sum_{y\in Y} f(g(y)).
\]

\smallskip
\noindent
\hyperref[prf:substitution-finite-set-summation]{\textit{Go to proof.}}
\end{proposition}
\end{propositionbox}

\begin{remark*}[Standard quantified statement]
\[
  g:Y\to X\text{ is a bijection}
  \Longrightarrow
  \sum_{x\in X}f(x)=\sum_{y\in Y}f(g(y)).
\]
\end{remark*}

\begin{remark*}[Interpretation]
A finite summation over \(X\) may be computed by summing over any set
bijective with \(X\), after composing the summand with the bijection.
\end{remark*}

\begin{dependencies}
\begin{itemize}
  \item \hyperref[def:summation-over-finite-set]{Summation over a Finite Set}
  \item \hyperref[prop:finite-summations-well-defined]{Finite Summations Are Well-Defined}
\end{itemize}
\end{dependencies}

\begin{propositionbox}{Proposition (Integer Interval as a Finite Set)}
% ---------------------------------------------------------
% Proposition: Integer Interval as a Finite Set
% Label: prop:integer-interval-finite-set-summation
% Source: Tao, Analysis I, Proposition 7.1.11(d)
% ---------------------------------------------------------
\begin{proposition}[Integer Interval as a Finite Set]
\label{prop:integer-interval-finite-set-summation}

Let \(n \le m\) be integers, and let
\[
  X:=\{i\in\mathbb{Z}:n\le i\le m\}.
\]
If \(a_i\) is a real number assigned to each integer \(i\in X\), then
\[
  \sum_{i=n}^{m} a_i
  =
  \sum_{i\in X} a_i.
\]

\smallskip
\noindent
\hyperref[prf:integer-interval-finite-set-summation]{\textit{Go to proof.}}
\end{proposition}
\end{propositionbox}

\begin{remark*}[Standard quantified statement]
\[
  X=\{i\in\mathbb Z:n\le i\le m\}
  \Longrightarrow
  \sum_{i=n}^{m}a_i=\sum_{i\in X}a_i.
\]
\end{remark*}

\begin{remark*}[Interpretation]
Summing over a consecutive integer interval agrees with summing over that
same interval regarded as a finite set.
\end{remark*}

\begin{dependencies}
\begin{itemize}
  \item \hyperref[def:finite-sum]{Finite Sum}
  \item \hyperref[def:summation-over-finite-set]{Summation over a Finite Set}
\end{itemize}
\end{dependencies}

\begin{propositionbox}{Proposition (Disjoint Union Finite Set Summation)}
% ---------------------------------------------------------
% Proposition: Disjoint Union Finite Set Summation
% Label: prop:disjoint-union-finite-set-summation
% Source: Tao, Analysis I, Proposition 7.1.11(e)
% ---------------------------------------------------------
\begin{proposition}[Disjoint Union Finite Set Summation]
\label{prop:disjoint-union-finite-set-summation}

Let \(X\) and \(Y\) be disjoint finite sets, so \(X\cap Y=\emptyset\),
and let \(f:X\cup Y\to\mathbb{R}\) be a function. Then
\[
  \sum_{z\in X\cup Y} f(z)
  =
  \left(\sum_{x\in X} f(x)\right)
  +
  \left(\sum_{y\in Y} f(y)\right).
\]

\smallskip
\noindent
\hyperref[prf:disjoint-union-finite-set-summation]{\textit{Go to proof.}}
\end{proposition}
\end{propositionbox}

\begin{remark*}[Standard quantified statement]
\[
  X\cap Y=\emptyset
  \Longrightarrow
  \sum_{z\in X\cup Y}f(z)
  =
  \sum_{x\in X}f(x)+\sum_{y\in Y}f(y).
\]
\end{remark*}

\begin{remark*}[Interpretation]
A sum over a disjoint union splits into the sum over each component set.
\end{remark*}

\begin{dependencies}
\begin{itemize}
  \item \hyperref[def:summation-over-finite-set]{Summation over a Finite Set}
\end{itemize}
\end{dependencies}

\begin{propositionbox}{Proposition (Addition Rule for Finite Set Summation)}
% ---------------------------------------------------------
% Proposition: Addition Rule for Finite Set Summation
% Label: prop:addition-rule-finite-set-summation
% Source: Tao, Analysis I, Proposition 7.1.11(f)
% ---------------------------------------------------------
\begin{proposition}[Addition Rule for Finite Set Summation]
\label{prop:addition-rule-finite-set-summation}

Let \(X\) be a finite set, and let \(f:X\to\mathbb{R}\) and
\(g:X\to\mathbb{R}\) be functions. Then
\[
  \sum_{x\in X} (f(x)+g(x))
  =
  \sum_{x\in X} f(x)
  +
  \sum_{x\in X} g(x).
\]

\smallskip
\noindent
\hyperref[prf:addition-rule-finite-set-summation]{\textit{Go to proof.}}
\end{proposition}
\end{propositionbox}

\begin{remark*}[Standard quantified statement]
\[
  \sum_{x\in X}(f(x)+g(x))
  =
  \sum_{x\in X}f(x)+\sum_{x\in X}g(x).
\]
\end{remark*}

\begin{remark*}[Interpretation]
Finite set summation distributes over pointwise addition.
\end{remark*}

\begin{dependencies}
\begin{itemize}
  \item \hyperref[def:summation-over-finite-set]{Summation over a Finite Set}
\end{itemize}
\end{dependencies}

\begin{propositionbox}{Proposition (Scalar Multiplication Rule for Finite Set Summation)}
% ---------------------------------------------------------
% Proposition: Scalar Multiplication Rule for Finite Set Summation
% Label: prop:scalar-multiplication-rule-finite-set-summation
% Source: Tao, Analysis I, Proposition 7.1.11(g)
% ---------------------------------------------------------
\begin{proposition}[Scalar Multiplication Rule for Finite Set Summation]
\label{prop:scalar-multiplication-rule-finite-set-summation}

Let \(X\) be a finite set, let \(f:X\to\mathbb{R}\) be a function, and
let \(c\) be a real number. Then
\[
  \sum_{x\in X} c f(x)
  =
  c \sum_{x\in X} f(x).
\]

\smallskip
\noindent
\hyperref[prf:scalar-multiplication-rule-finite-set-summation]{\textit{Go to proof.}}
\end{proposition}
\end{propositionbox}

\begin{remark*}[Standard quantified statement]
\[
  \sum_{x\in X}c f(x)=c\sum_{x\in X}f(x).
\]
\end{remark*}

\begin{remark*}[Interpretation]
A constant scalar may be pulled outside a finite set summation.
\end{remark*}

\begin{dependencies}
\begin{itemize}
  \item \hyperref[def:summation-over-finite-set]{Summation over a Finite Set}
\end{itemize}
\end{dependencies}

\begin{propositionbox}{Proposition (Monotonicity for Finite Set Summation)}
% ---------------------------------------------------------
% Proposition: Monotonicity for Finite Set Summation
% Label: prop:monotonicity-finite-set-summation
% Source: Tao, Analysis I, Proposition 7.1.11(h)
% ---------------------------------------------------------
\begin{proposition}[Monotonicity for Finite Set Summation]
\label{prop:monotonicity-finite-set-summation}

Let \(X\) be a finite set, and let \(f:X\to\mathbb{R}\) and
\(g:X\to\mathbb{R}\) be functions such that \(f(x)\le g(x)\) for all
\(x\in X\). Then
\[
  \sum_{x\in X} f(x)
  \le
  \sum_{x\in X} g(x).
\]

\smallskip
\noindent
\hyperref[prf:monotonicity-finite-set-summation]{\textit{Go to proof.}}
\end{proposition}
\end{propositionbox}

\begin{remark*}[Standard quantified statement]
\[
  \bigl(\forall x\in X,\ f(x)\le g(x)\bigr)
  \Longrightarrow
  \sum_{x\in X} f(x)
  \le
  \sum_{x\in X} g(x).
\]
\end{remark*}

\begin{remark*}[Interpretation]
Pointwise order between two functions on a finite set gives the same
order between their finite sums.
\end{remark*}

\begin{dependencies}
\begin{itemize}
  \item \hyperref[def:summation-over-finite-set]{Summation over a Finite Set}
\end{itemize}
\end{dependencies}

\begin{propositionbox}{Proposition (Triangle Inequality for Finite Set Summation)}
% ---------------------------------------------------------
% Proposition: Triangle Inequality for Finite Set Summation
% Label: prop:triangle-inequality-finite-set-summation
% Source: Tao, Analysis I, Proposition 7.1.11(i)
% ---------------------------------------------------------
\begin{proposition}[Triangle Inequality for Finite Set Summation]
\label{prop:triangle-inequality-finite-set-summation}

Let \(X\) be a finite set, and let \(f:X\to\mathbb{R}\) be a function.
Then
\[
  \left\lvert \sum_{x\in X} f(x) \right\rvert
  \le
  \sum_{x\in X} \lvert f(x)\rvert.
\]

\smallskip
\noindent
\hyperref[prf:triangle-inequality-finite-set-summation]{\textit{Go to proof.}}
\end{proposition}
\end{propositionbox}

\begin{remark*}[Standard quantified statement]
\[
  \left\lvert \sum_{x\in X}f(x)\right\rvert
  \le
  \sum_{x\in X}|f(x)|.
\]
\end{remark*}

\begin{remark*}[Interpretation]
The absolute value of a finite set summation is bounded by the
corresponding finite set summation of absolute values.
\end{remark*}

\begin{dependencies}
\begin{itemize}
  \item \hyperref[def:summation-over-finite-set]{Summation over a Finite Set}
\end{itemize}
\end{dependencies}

\begin{lemmabox}{Lemma (Iterated Summation over a Product)}
% ---------------------------------------------------------
% Lemma: Iterated Summation over a Product
% Label: lem:iterated-summation-over-product
% Source: Tao, Analysis I, Lemma 7.1.13
% ---------------------------------------------------------
\begin{lemma}[Iterated Summation over a Product]
\label{lem:iterated-summation-over-product}

Let \(X\) and \(Y\) be finite sets, and let
\(f:X\times Y\to\mathbb{R}\) be a function. Then
\[
  \sum_{x\in X}
  \left(
    \sum_{y\in Y} f(x,y)
  \right)
  =
  \sum_{(x,y)\in X\times Y} f(x,y).
\]

\smallskip
\noindent
\hyperref[prf:iterated-summation-over-product]{\textit{Go to proof.}}
\end{lemma}
\end{lemmabox}

\begin{remark*}[Standard quantified statement]
\[
  \sum_{x\in X}\left(\sum_{y\in Y}f(x,y)\right)
  =
  \sum_{(x,y)\in X\times Y}f(x,y).
\]
\end{remark*}

\begin{remark*}[Interpretation]
Summing first over \(Y\) for each \(x\in X\), and then summing those
values over \(X\), gives the same value as summing over the product set
\(X\times Y\).
\end{remark*}

\begin{dependencies}
\begin{itemize}
  \item \hyperref[def:summation-over-finite-set]{Summation over a Finite Set}
\end{itemize}
\end{dependencies}

\begin{corollarybox}{Corollary (Fubini's Theorem for Finite Series)}
% ---------------------------------------------------------
% Corollary: Fubini's Theorem for Finite Series
% Label: cor:fubini-theorem-finite-series
% Source: Tao, Analysis I, Corollary 7.1.14
% ---------------------------------------------------------
\begin{corollary}[Fubini's Theorem for Finite Series]
\label{cor:fubini-theorem-finite-series}

Let \(X\) and \(Y\) be finite sets, and let
\(f:X\times Y\to\mathbb{R}\) be a function. Then
\[
  \sum_{x\in X}
  \left(
    \sum_{y\in Y} f(x,y)
  \right)
  =
  \sum_{(x,y)\in X\times Y} f(x,y)
  =
  \sum_{(y,x)\in Y\times X} f(x,y)
  =
  \sum_{y\in Y}
  \left(
    \sum_{x\in X} f(x,y)
  \right).
\]

\smallskip
\noindent
\hyperref[prf:fubini-theorem-finite-series]{\textit{Go to proof.}}
\end{corollary}
\end{corollarybox}

\begin{remark*}[Standard quantified statement]
\[
  \sum_{x\in X}\left(\sum_{y\in Y} f(x,y)\right)
  =
  \sum_{y\in Y}\left(\sum_{x\in X} f(x,y)\right).
\]
\end{remark*}

\begin{remark*}[Interpretation]
For a function on a finite product set, the order of summation over the two
finite factors may be interchanged.
\end{remark*}

\begin{dependencies}
\begin{itemize}
  \item \hyperref[lem:iterated-summation-over-product]{Iterated Summation over a Product}
  \item \hyperref[def:summation-over-finite-set]{Summation over a Finite Set}
\end{itemize}
\end{dependencies}

\begin{lemmabox}{Lemma (Reindexing a Finite Sum)}
% ---------------------------------------------------------
% Lemma: Reindexing a Finite Sum
% Label: lem:finite-sum-reindexing
% Source: Tao, Analysis I, Lemma 7.1.4(b)
% ---------------------------------------------------------
\begin{lemma}[Reindexing a Finite Sum]
\label{lem:finite-sum-reindexing}

Let \(m \le n\) be integers, let \(k\) be an integer, and let \(a_i\)
be a real number assigned to each integer \(i\) satisfying \(m \le i
\le n\). Then
\[
  \sum_{i=m}^{n} a_i
  =
  \sum_{j=m+k}^{n+k} a_{j-k}.
\]

\smallskip
\noindent
\hyperref[prf:finite-sum-reindexing]{\textit{Go to proof.}}
\end{lemma}
\end{lemmabox}

\begin{remark*}[Standard quantified statement]
\[
  \sum_{i=m}^{n}a_i
  =
  \sum_{j=m+k}^{n+k}a_{j-k}.
\]
\end{remark*}

\begin{remark*}[Interpretation]
Shifting every index by the same integer changes the labels of the
summation but not the ordered list of terms being summed.
\end{remark*}

\begin{dependencies}
\begin{itemize}
  \item \hyperref[def:recursive-finite-sum]{Recursive Finite Sum}
\end{itemize}
\end{dependencies}

\begin{lemmabox}{Lemma (Finite Sum of Sums)}
% ---------------------------------------------------------
% Lemma: Finite Sum of Sums
% Label: lem:finite-sum-of-sums
% Source: Tao, Analysis I, Lemma 7.1.4(c)
% ---------------------------------------------------------
\begin{lemma}[Finite Sum of Sums]
\label{lem:finite-sum-of-sums}

Let \(m \le n\) be integers, and let \(a_i\) and \(b_i\) be real
numbers assigned to each integer \(i\) satisfying \(m \le i \le n\).
Then
\[
  \sum_{i=m}^{n} (a_i+b_i)
  =
  \left(\sum_{i=m}^{n} a_i\right)
  +
  \left(\sum_{i=m}^{n} b_i\right).
\]

\smallskip
\noindent
\hyperref[prf:finite-sum-of-sums]{\textit{Go to proof.}}
\end{lemma}
\end{lemmabox}

\begin{remark*}[Standard quantified statement]
\[
  \sum_{i=m}^{n}(a_i+b_i)
  =
  \sum_{i=m}^{n}a_i+\sum_{i=m}^{n}b_i.
\]
\end{remark*}

\begin{remark*}[Interpretation]
Finite summation distributes over termwise addition.
\end{remark*}

\begin{dependencies}
\begin{itemize}
  \item \hyperref[def:recursive-finite-sum]{Recursive Finite Sum}
\end{itemize}
\end{dependencies}

\begin{lemmabox}{Lemma (Finite Sum of Scalar Multiples)}
% ---------------------------------------------------------
% Lemma: Finite Sum of Scalar Multiples
% Label: lem:finite-sum-scalar-multiple
% Source: Tao, Analysis I, Lemma 7.1.4(d)
% ---------------------------------------------------------
\begin{lemma}[Finite Sum of Scalar Multiples]
\label{lem:finite-sum-scalar-multiple}

Let \(m \le n\) be integers, let \(a_i\) be a real number assigned to
each integer \(i\) satisfying \(m \le i \le n\), and let
\(c \in \mathbb{R}\). Then
\[
  \sum_{i=m}^{n} (c a_i)
  =
  c\left(\sum_{i=m}^{n} a_i\right).
\]

\smallskip
\noindent
\hyperref[prf:finite-sum-scalar-multiple]{\textit{Go to proof.}}
\end{lemma}
\end{lemmabox}

\begin{remark*}[Standard quantified statement]
\[
  \sum_{i=m}^{n}(c a_i)
  =
  c\sum_{i=m}^{n}a_i.
\]
\end{remark*}

\begin{remark*}[Interpretation]
A common scalar factor may be pulled outside a finite sum.
\end{remark*}

\begin{dependencies}
\begin{itemize}
  \item \hyperref[def:recursive-finite-sum]{Recursive Finite Sum}
\end{itemize}
\end{dependencies}

\begin{lemmabox}{Lemma (Triangle Inequality for Finite Series)}
% ---------------------------------------------------------
% Lemma: Triangle Inequality for Finite Series
% Label: lem:triangle-inequality-finite-series
% Source: Tao, Analysis I, Lemma 7.1.4(e)
% ---------------------------------------------------------
\begin{lemma}[Triangle Inequality for Finite Series]
\label{lem:triangle-inequality-finite-series}

Let \(m \le n\) be integers, and let \(a_i\) be a real number assigned
to each integer \(i\) satisfying \(m \le i \le n\). Then
\[
  \left\lvert \sum_{i=m}^{n} a_i \right\rvert
  \le
  \sum_{i=m}^{n} \lvert a_i\rvert.
\]

\smallskip
\noindent
\hyperref[prf:triangle-inequality-finite-series]{\textit{Go to proof.}}
\end{lemma}
\end{lemmabox}

\begin{remark*}[Standard quantified statement]
\[
  \forall m,n \in \mathbb{Z},\quad
  m \le n
  \Longrightarrow
  \left\lvert \sum_{i=m}^{n} a_i \right\rvert
  \le
  \sum_{i=m}^{n} \lvert a_i\rvert.
\]
\end{remark*}

\begin{remark*}[Predicate reading]
\[
  m\le n
  \Longrightarrow
  \left|\sum_{i=m}^{n}a_i\right|
  \le
  \sum_{i=m}^{n}|a_i|.
\]
\end{remark*}

\begin{remark*}[Interpretation]
The absolute value of a finite series is bounded by the corresponding
finite series of absolute values.
\end{remark*}

\begin{dependencies}
\begin{itemize}
  \item \hyperref[def:recursive-finite-sum]{Recursive Finite Sum}
\end{itemize}
\end{dependencies}

\begin{lemmabox}{Lemma (Comparison Test for Finite Series)}
% ---------------------------------------------------------
% Lemma: Comparison Test for Finite Series
% Label: lem:comparison-test-finite-series
% Source: Tao, Analysis I, Lemma 7.1.4(f)
% ---------------------------------------------------------
\begin{lemma}[Comparison Test for Finite Series]
\label{lem:comparison-test-finite-series}

Let \(m \le n\) be integers, and let \(a_i\) and \(b_i\) be real
numbers assigned to each integer \(i\) satisfying \(m \le i \le n\).
Suppose that \(a_i \le b_i\) for all \(i\) satisfying \(m \le i \le n\).
Then
\[
  \sum_{i=m}^{n} a_i
  \le
  \sum_{i=m}^{n} b_i.
\]

\smallskip
\noindent
\hyperref[prf:comparison-test-finite-series]{\textit{Go to proof.}}
\end{lemma}
\end{lemmabox}

\begin{remark*}[Standard quantified statement]
\[
  \forall i\,(m \le i \le n \Longrightarrow a_i \le b_i)
  \Longrightarrow
  \sum_{i=m}^{n} a_i
  \le
  \sum_{i=m}^{n} b_i.
\]
\end{remark*}

\begin{remark*}[Predicate reading]
\[
  \left(\forall i,\ m\le i\le n\Rightarrow a_i\le b_i\right)
  \Longrightarrow
  \sum_{i=m}^{n}a_i
  \le
  \sum_{i=m}^{n}b_i.
\]
\end{remark*}

\begin{remark*}[Interpretation]
Termwise comparison of two finite series gives comparison of their finite
sums.
\end{remark*}

\begin{dependencies}
\begin{itemize}
  \item \hyperref[def:recursive-finite-sum]{Recursive Finite Sum}
\end{itemize}
\end{dependencies}
